\documentclass[11pt,a4paper]{article}
\usepackage[utf8]{inputenc}
\usepackage[T1]{fontenc}
\usepackage[french]{babel}
\usepackage{amsmath,amssymb}
\usepackage{booktabs}
\usepackage{siunitx}
\usepackage{geometry}
\usepackage{enumitem}
\usepackage{xcolor}
\usepackage{tikz}
\usepackage{pgfplots}
\usepackage{float}
\usepackage{array}
\usepackage{tabularx}
\usepackage{colortbl}
\usepackage{fancyhdr}

\pgfplotsset{compat=1.18}
\usetikzlibrary{shapes.geometric, positioning, calc, patterns, arrows.meta}

\geometry{margin=2cm, headheight=14pt}
\sisetup{output-decimal-marker={,},group-separator={\,}}

% ------- Couleurs ---------------------------------------------
\definecolor{calchead}{RGB}{230,230,230}
\definecolor{calcline}{RGB}{150,150,150}
\definecolor{calcform}{RGB}{0,102,204}
\definecolor{calcval}{RGB}{0,0,0}
\definecolor{boitecol}{RGB}{173,216,230}
\definecolor{medcol}{RGB}{200,60,60}
\definecolor{encadre}{RGB}{255,248,220}

% ------- Environnement Calc ------------------------------------
\newenvironment{calcsheet}[1][\footnotesize]
  {\begin{center}#1
   \begin{tabular}{|>{\columncolor{calchead}}c|l|l|l|l|}
   \hline
   \rowcolor{calchead}
   & \textbf{A} & \textbf{B} & \textbf{C} & \textbf{D} \\ \hline}
  {\end{tabular}\end{center}}

\newcommand{\fc}[1]{\textcolor{calcform}{\texttt{#1}}}

% ------- Encadré énoncé (version corrigée) --------------------
\newsavebox{\enoncebox}
\newenvironment{enonce}
  {\par\medskip\begin{center}%
   \begin{lrbox}{\enoncebox}%
   \begin{minipage}{0.9\textwidth}%
   \textbf{Énoncé.}\ \itshape\par\smallskip}
  {\end{minipage}%
   \end{lrbox}%
   \setlength{\fboxsep}{10pt}%
   \colorbox{encadre}{\usebox{\enoncebox}}%
   \end{center}\par\medskip}

% ------- Pagination -------------------------------------------
\pagestyle{fancy}
\fancyhf{}
\rhead{\small Statistique descriptive}
\lhead{\small STS BIOAL}
\cfoot{\thepage}

\title{\textbf{Correction détaillée — Statistique descriptive}\\[0.3em]
\large Calculs manuels, formules tableurs \& simulations Calc}
\author{STS BIOAL — Module 5}
\date{}

\begin{document}
\maketitle

\tableofcontents
\newpage

% ============================================================
\section{Rappel de l'énoncé}

\subsection*{Contexte général}

Dans le cadre du module de statistique descriptive (STS Bioanalyse),
un laboratoire d'analyses médicales contrôle la concentration en
glucose d'un lot de sérums. La valeur cible du fabricant est fixée
à \( 15{,}5 \)~g/L. Le technicien doit vérifier la conformité du lot
et diagnostiquer d'éventuelles anomalies.

\subsection*{Énoncé de la partie « Bonus »}

\begin{enonce}
On dispose de \( n = 40 \) mesures de concentration en glucose
(en g/L) réalisées sur un même lot :

\smallskip
\noindent
12,1 \quad 12,6 \quad 13,0 \quad 13,3 \quad 13,6 \quad 13,9 \quad 14,0 \quad 14,2 \quad 14,4 \quad 14,5 \\
14,7 \quad 14,9 \quad 15,0 \quad 15,1 \quad 15,2 \quad 15,3 \quad 15,4 \quad 15,5 \quad 15,6 \quad 15,8 \\
15,9 \quad 16,0 \quad 16,1 \quad 16,3 \quad 16,4 \quad 16,5 \quad 16,6 \quad 16,8 \quad 16,9 \quad 17,0 \\
17,2 \quad 17,4 \quad 17,5 \quad 17,7 \quad 17,9 \quad 18,1 \quad 18,4 \quad 18,6 \quad 18,9 \quad 19,1

\smallskip
\noindent\textbf{Questions bonus :}
\begin{enumerate}[label=\textbf{C.\arabic*.}, leftmargin=*]
  \item Calculer la moyenne exacte et l'écart-type exact à partir
        des valeurs brutes, puis comparer avec la méthode des classes.
  \item Démontrer que la somme des écarts à la moyenne est nulle,
        et vérifier numériquement.
  \item Démontrer l'identité de König
        \( \sum (x_i-\bar{x})^2 = \sum x_i^2 - n\bar{x}^2 \)
        et en déduire la variance exacte.
  \item Quelle proportion des échantillons dépasse
        \( \bar{x} + \sigma \) ?
  \item La valeur \( 19{,}1 \) est-elle aberrante selon la règle de Tukey ?
  \item Que deviennent moyenne, variance, écart-type et coefficient
        de variation si l'on convertit les concentrations en mg/dL
        (multiplication par 100) ?
  \item Établir la relation entre les coefficients de la régression
        de \( y \) en \( x \) et ceux de la régression inverse de
        \( x \) en \( y \).
  \item Rédiger une conclusion biologique argumentée sur la conformité
        du lot.
\end{enumerate}
\end{enonce}

% ============================================================
\section{Notation \(\sum\) : signification et règles de calcul}

\subsection{Que signifie le symbole \(\sum\) ?}

Le symbole \(\sum\) (lettre grecque majuscule \emph{sigma}) est
l'opérateur de \textbf{sommation}. Il indique que l'on doit
\emph{additionner} une quantité pour toutes les valeurs d'un
indice variant entre deux bornes.

\medskip
\noindent Formellement :
\[
\sum_{i=1}^{n} x_i
\;=\;
x_1 + x_2 + x_3 + \cdots + x_{n-1} + x_n.
\]

\noindent Décomposition des éléments :
\begin{center}
\begin{tabular}{|c|l|}
\hline
\rowcolor{calchead}
\textbf{Élément} & \textbf{Rôle} \\ \hline
\( \sum \)        & « somme de \dots » \\ \hline
\( i = 1 \)       & \emph{indice de départ} (on commence à \( i = 1 \)) \\ \hline
\( n \)           & \emph{indice d'arrivée} (on s'arrête à \( i = n \)) \\ \hline
\( x_i \)         & \emph{terme général} à additionner \\ \hline
\end{tabular}
\end{center}

\subsection{Exemple concret sur nos données}

Pour la série des 40 concentrations, l'écriture
\[
\sum_{i=1}^{40} x_i
\]
signifie littéralement :
\[
x_1 + x_2 + x_3 + \cdots + x_{39} + x_{40}
\]
\[
= 12{,}1 + 12{,}6 + 13{,}0 + \cdots + 18{,}9 + 19{,}1 = 633{,}4.
\]

De même,
\[
\sum_{i=1}^{40} x_i^2 = x_1^2 + x_2^2 + \cdots + x_{40}^2
= 12{,}1^2 + 12{,}6^2 + \cdots + 19{,}1^2 = 10\,151{,}16.
\]

\subsection{Règles de calcul à connaître}

\begin{enumerate}[label=\textbf{R\arabic*.}, leftmargin=*]
  \item \textbf{Somme d'une constante.} Si \( c \) est une constante
        (indépendante de \( i \)), alors :
        \[
        \sum_{i=1}^{n} c = n \times c.
        \]
        \emph{Exemple :} \( \sum_{i=1}^{40} 3 = 40 \times 3 = 120 \).

  \item \textbf{Linéarité.} Pour deux suites \( (x_i) \) et \( (y_i) \)
        et deux constantes \( a, b \) :
        \[
        \sum_{i=1}^{n} (a\,x_i + b\,y_i) = a \sum_{i=1}^{n} x_i
        + b \sum_{i=1}^{n} y_i.
        \]

  \item \textbf{Factorisation d'une constante.}
        \[
        \sum_{i=1}^{n} (c \cdot x_i) = c \sum_{i=1}^{n} x_i.
        \]

  \item \textbf{Attention !} En général :
        \[
        \sum_{i=1}^{n} x_i^2 \neq \left(\sum_{i=1}^{n} x_i\right)^{2}.
        \]
        Dans notre cas : \( \sum x_i^2 = 10\,151{,}16 \)
        tandis que \( \left(\sum x_i\right)^2 = 633{,}4^2 \approx 401\,195 \).
        Ces deux quantités sont totalement différentes !
\end{enumerate}

\subsection{Notations abrégées utilisées dans la suite}

Par souci de lisibilité, on notera souvent :
\[
\sum x_i \quad\text{pour}\quad \sum_{i=1}^{n} x_i,
\qquad
\sum x_i^2 \quad\text{pour}\quad \sum_{i=1}^{n} x_i^2.
\]

\noindent De même, pour la somme des carrés des écarts :
\[
\sum (x_i - \bar{x})^2
\quad\text{signifie}\quad
\sum_{i=1}^{n} (x_i - \bar{x})^2
= (x_1-\bar{x})^2 + (x_2-\bar{x})^2 + \cdots + (x_n-\bar{x})^2.
\]

\subsection{Traduction dans le tableur}

Dans LibreOffice Calc ou Excel, la somme \( \sum x_i \) se traduit
directement par la fonction :

\begin{center}
\fc{=SOMME(A2:A41)} \quad (français) \qquad
\fc{=SUM(A2:A41)} \quad (anglais)
\end{center}

\noindent Pour la somme des carrés \( \sum x_i^2 \) :
\begin{center}
\fc{=SOMME.SQ(A2:A41)} \quad (français) \qquad
\fc{=SUMSQ(A2:A41)} \quad (anglais)
\end{center}

\noindent Ces fonctions réalisent \emph{automatiquement} la
sommation décrite par le symbole \( \sum \) : elles parcourent
toutes les cellules de la plage \( A2:A41 \) et additionnent
leurs valeurs.

% ============================================================
\section{Mise en place des données}

\subsection{Série brute}

Un laboratoire de bioanalyse contrôle la concentration en glucose
(\si{g/L}) d'un lot de \( n = 40 \) échantillons sériques :

\begin{center}
\footnotesize
\begin{tabular}{|l|l|l|l|l|l|l|l|l|l|}
\hline
\rowcolor{calchead}
12,1 & 12,6 & 13,0 & 13,3 & 13,6 & 13,9 & 14,0 & 14,2 & 14,4 & 14,5 \\
14,7 & 14,9 & 15,0 & 15,1 & 15,2 & 15,3 & 15,4 & 15,5 & 15,6 & 15,8 \\
15,9 & 16,0 & 16,1 & 16,3 & 16,4 & 16,5 & 16,6 & 16,8 & 16,9 & 17,0 \\
17,2 & 17,4 & 17,5 & 17,7 & 17,9 & 18,1 & 18,4 & 18,6 & 18,9 & 19,1 \\
\hline
\end{tabular}
\end{center}

\medskip
\noindent Sommes fondamentales (calculs vérifiés) :
\[
\sum_{i=1}^{40} x_i = 633{,}4
\qquad
\sum_{i=1}^{40} x_i^{2} = 10\,151{,}16
\]

\subsection{Simulation de la feuille Calc}

On saisit les valeurs dans la colonne \( A \) :

\begin{calcsheet}
1 & Concentration (g/L) & & & \\ \hline
2 & 12,1 & & & \\ \hline
3 & 12,6 & & & \\ \hline
4 & 13,0 & & & \\ \hline
\(\vdots\) & \(\vdots\) & & & \\ \hline
41 & 19,1 & & & \\ \hline
42 & & & & \\ \hline
43 & \( n \)  & \fc{=NB(A2:A41)} & 40 & \\ \hline
44 & \( \bar{x} \) & \fc{=MOYENNE(A2:A41)} & 15,835 & \\ \hline
45 & \( \sum x_i \) & \fc{=SOMME(A2:A41)} & 633,4 & \\ \hline
46 & \( \sum x_i^2 \) & \fc{=SOMME.SQ(A2:A41)} & 10151,16 & \\ \hline
\end{calcsheet}

\noindent\textbf{Équivalents anglais (Excel) :}
\fc{COUNT}, \fc{AVERAGE}, \fc{SUM}, \fc{SUMSQ}.

% ============================================================
\section{Indicateurs de position}

\subsection{Moyenne arithmétique}

\[
\bar{x} = \frac{1}{n}\sum_{i=1}^{n} x_i
= \frac{633{,}4}{40}
= \boxed{15{,}835\ \si{g/L}}.
\]

\noindent La formule \( \bar{x} = \tfrac{1}{n}\sum x_i \) se lit
« \emph{la moyenne est égale à la somme de toutes les valeurs,
divisée par leur nombre} ».

\subsection{Médiane}

Série déjà triée, \( n = 40 \) (pair) : \( Me = \frac{x_{20}+x_{21}}{2} \).
\[
x_{20} = 15{,}8,\quad x_{21} = 15{,}9
\quad\Longrightarrow\quad
Me = \boxed{15{,}85\ \si{g/L}}.
\]

\subsection{Quartiles}

Position : \( p_k = \frac{k(n+1)}{4} \).

\medskip
\noindent Pour \( Q_1 \) : \( p_1 = 10{,}25 \), donc
\[
Q_1 = x_{10} + 0{,}25\,(x_{11}-x_{10}) = 14{,}5 + 0{,}25 \times 0{,}2
= \boxed{14{,}55}.
\]

\noindent Pour \( Q_3 \) : \( p_3 = 30{,}75 \), donc
\[
Q_3 = x_{30} + 0{,}75\,(x_{31}-x_{30}) = 17{,}0 + 0{,}75 \times 0{,}2
= \boxed{17{,}15}.
\]

\[
IQR = Q_3 - Q_1 = 17{,}15 - 14{,}55 = \boxed{2{,}60\ \si{g/L}}.
\]

\begin{calcsheet}
48 & Médiane \( Me \) & \fc{=MEDIANE(A2:A41)} & 15,85 & \\ \hline
49 & \( Q_1 \) & \fc{=QUARTILE.INC(A2:A41;1)} & 14,55 & \\ \hline
50 & \( Q_3 \) & \fc{=QUARTILE.INC(A2:A41;3)} & 17,15 & \\ \hline
51 & \( IQR \) & \fc{=D50-D49} & 2,60 & \\ \hline
\end{calcsheet}

\noindent\textbf{En anglais} : \fc{=MEDIAN(...)}, \fc{=QUARTILE.INC(...,k)}.

% ============================================================
\section{Indicateurs de dispersion}

\subsection{Variance et écart-type (formule de König)}

La formule de König permet de calculer la somme des carrés des
écarts (SCE) sans avoir à calculer chaque écart individuellement :
\[
\sum_{i=1}^{n}(x_i-\bar{x})^2 = \sum_{i=1}^{n} x_i^2 - n\bar{x}^2.
\]

\noindent\textbf{Signification terme à terme :}
\begin{itemize}
  \item \( \sum x_i^2 \) : somme des carrés de toutes les valeurs
        (calculée directement sur les données brutes) ;
  \item \( n\bar{x}^2 \) : produit de l'effectif \( n \) par le carré
        de la moyenne.
\end{itemize}

\[
V = \frac{1}{n}\left[\sum x_i^2 - n\bar{x}^2\right]
= \frac{10\,151{,}16 - 40\times 15{,}835^2}{40}
\]

\[
n\bar{x}^2 = 40 \times 250{,}747225 = 10\,029{,}889
\]
\[
SCE = 10\,151{,}16 - 10\,029{,}889 = 121{,}271
\]
\[
V = \frac{121{,}271}{40} = \boxed{3{,}0318}
\qquad
\sigma = \sqrt{3{,}0318} \approx \boxed{1{,}7412\ \si{g/L}}
\]

\noindent Écart-type échantillonnal :
\[
s^2 = \frac{SCE}{n-1} = \frac{121{,}271}{39} \approx 3{,}1095
\qquad
s \approx \boxed{1{,}7634\ \si{g/L}}
\]

\begin{calcsheet}
53 & SCE & \fc{=SOMME.SQ(A2:A41)-A43*A44\^{}2} & 121,271 & \\ \hline
54 & \( V \) (pop.) & \fc{=VAR.P(A2:A41)} & 3,0318 & \\ \hline
55 & \( \sigma \) (pop.) & \fc{=ECARTYPE.P(A2:A41)} & 1,7412 & \\ \hline
56 & \( s^2 \) (éch.) & \fc{=VAR.S(A2:A41)} & 3,1095 & \\ \hline
57 & \( s \) (éch.) & \fc{=ECARTYPE.S(A2:A41)} & 1,7634 & \\ \hline
58 & Étendue & \fc{=MAX(A2:A41)-MIN(A2:A41)} & 7,0 & \\ \hline
59 & \( CV \) & \fc{=ECARTYPE.P(A2:A41)/MOYENNE(A2:A41)} & 0,1100 & \\ \hline
\end{calcsheet}

% ============================================================
\section{Histogramme des effectifs par classes}

\subsection{Construction du tableau de classes}

\begin{center}
\small
\begin{tabular}{|c|c|c|c|}
\hline
\rowcolor{calchead}
\textbf{Classe} & \textbf{Centre \( c_i \)} & \textbf{Effectif \( n_i \)} & \textbf{Fréquence} \\ \hline
\( [12\,;14[ \) & 13 & 4  & 0,100 \\ \hline
\( [14\,;16[ \) & 15 & 17 & 0,425 \\ \hline
\( [16\,;18[ \) & 17 & 14 & 0,350 \\ \hline
\( [18\,;20[ \) & 19 & 5  & 0,125 \\ \hline
\textbf{Total} & --- & \textbf{40} & \textbf{1,000} \\ \hline
\end{tabular}
\end{center}

\subsection{Histogramme}

\begin{figure}[H]
\centering
\begin{tikzpicture}
\begin{axis}[
  width=12cm, height=7cm,
  ybar interval,
  xlabel={Concentration (\si{g/L})},
  ylabel={Effectif \( n_i \)},
  xmin=11, xmax=21,
  ymin=0, ymax=20,
  xtick={12,14,16,18,20},
  ytick={0,5,10,15,20},
  grid=major,
  grid style={dashed,gray!40},
  bar width=2cm,
  area style,
  fill=boitecol!80,
  draw=blue!60!black,
  every axis plot/.append style={line width=1pt},
]
\addplot coordinates {(12,4) (14,17) (16,14) (18,5) (20,0)};
\end{axis}
\end{tikzpicture}
\caption{Histogramme des effectifs par classes de largeur \SI{2}{g/L}.}
\end{figure}

\subsection{Moyenne et variance par classes}

\[
\bar{x}_{\text{cl}} = \frac{\sum n_i c_i}{\sum n_i}
= \frac{4\times 13 + 17\times 15 + 14\times 17 + 5\times 19}{40}
= \frac{640}{40} = \boxed{16{,}0}
\]

\[
V_{\text{cl}} = \frac{4\cdot 9 + 17\cdot 1 + 14\cdot 1 + 5\cdot 9}{40}
= \frac{112}{40} = \boxed{2{,}8}
\qquad
\sigma_{\text{cl}} \approx \boxed{1{,}6733}.
\]

\begin{calcsheet}
61 & Centre \( c_i \) & Effectif \( n_i \) & & \\ \hline
62 & 13 & 4  & & \\ \hline
63 & 15 & 17 & & \\ \hline
64 & 17 & 14 & & \\ \hline
65 & 19 & 5  & & \\ \hline
66 & \( \bar{x}_{\text{cl}} \) & \fc{=SOMMEPROD(A62:A65;B62:B65)/SOMME(B62:B65)} & 16,0 & \\ \hline
67 & \( V_{\text{cl}} \) & \fc{=SOMMEPROD(B62:B65;(A62:A65-A66)\^{}2)/SOMME(B62:B65)} & 2,8 & \\ \hline
\end{calcsheet}

\noindent\textbf{Remarque :} la méthode des classes sous-estime la variance
(\( V_{\text{cl}} = 2{,}8 < V_{\text{exact}} = 3{,}03 \)),
car la dispersion interne aux classes est « perdue ».

% ============================================================
\section{Boîte à moustaches}

\begin{figure}[H]
\centering
\begin{tikzpicture}
\begin{axis}[
  width=13cm, height=5cm,
  xmin=10, xmax=22,
  ymin=-1.2, ymax=1.2,
  xtick={11,12,13,14,15,16,17,18,19,20,21},
  ytick=\empty,
  axis y line=none,
  axis x line=bottom,
  xlabel={Concentration (\si{g/L})},
  clip=false,
]

% --- Bornes de Tukey (pointillés) -----------------------------
\draw[dashed,gray]
  (axis cs:10.65,-1.1) -- (axis cs:10.65,1.0)
  node[above,black,font=\scriptsize]{10,65};
\draw[dashed,gray]
  (axis cs:21.05,-1.1) -- (axis cs:21.05,1.0)
  node[above,black,font=\scriptsize]{21,05};

% --- Boîte Q1 .. Q3 -------------------------------------------
\fill[boitecol!80]
  (axis cs:14.55,-0.4) rectangle (axis cs:17.15,0.4);
\draw[thick]
  (axis cs:14.55,-0.4) rectangle (axis cs:17.15,0.4);

% --- Médiane --------------------------------------------------
\draw[medcol, line width=2pt]
  (axis cs:15.85,-0.4) -- (axis cs:15.85,0.4);

% --- Moustaches -----------------------------------------------
\draw[thick] (axis cs:12.1,0) -- (axis cs:14.55,0);
\draw[thick] (axis cs:12.1,-0.25) -- (axis cs:12.1,0.25);
\draw[thick] (axis cs:17.15,0) -- (axis cs:19.1,0);
\draw[thick] (axis cs:19.1,-0.25) -- (axis cs:19.1,0.25);

% --- Étiquettes -----------------------------------------------
\node[below,font=\small] at (axis cs:15.85,-0.5)
  {\(Me=15{,}85\)};
\node[above,font=\small] at (axis cs:14.55,0.5) {\(Q_1\)};
\node[above,font=\small] at (axis cs:17.15,0.5) {\(Q_3\)};
\node[above,font=\small] at (axis cs:12.1,0.35) {min};
\node[above,font=\small] at (axis cs:19.1,0.35) {max};

% --- Titre ----------------------------------------------------
\node[font=\bfseries] at (axis cs:16,1.1)
  {Boîte à moustaches (règle de Tukey)};

\end{axis}
\end{tikzpicture}
\caption{Boîte à moustaches. La valeur maximale \( 19{,}1 \) reste
  en deçà de la borne supérieure \( 21{,}05 \) : aucune valeur n'est
  aberrante.}
\end{figure}

% ============================================================
\section{Détection d'une valeur aberrante}

Règle de Tukey : une valeur est aberrante si
\[
x < Q_1 - 1{,}5\,IQR \quad\text{ou}\quad x > Q_3 + 1{,}5\,IQR.
\]

\[
Q_1 - 1{,}5\,IQR = 14{,}55 - 3{,}90 = 10{,}65,
\qquad
Q_3 + 1{,}5\,IQR = 17{,}15 + 3{,}90 = 21{,}05.
\]

Intervalle admissible : \( [10{,}65\,;\,21{,}05] \).

\medskip
La valeur maximale \( 19{,}1 \) est \emph{strictement à l'intérieur} :
\( \boxed{\text{pas d'aberration}} \).

\begin{calcsheet}
69 & Borne inf. & \fc{=D49-1,5*D51} & 10,65 & \\ \hline
70 & Borne sup. & \fc{=D50+1,5*D51} & 21,05 & \\ \hline
71 & Aberrant ? & \fc{=SI(OU(A2<D69;A2>D70);"OUI";"NON")} & NON & \\ \hline
\end{calcsheet}

% ============================================================
\section{Régression linéaire — illustration}

On considère ici un exemple de droite \( y = ax + b \) avec les
indicateurs de qualité associés. La droite de régression de \( y \)
en \( x \) a pour équation :
\[
a = \frac{\mathrm{Cov}(x,y)}{V(x)},
\qquad
b = \bar{y} - a\bar{x},
\qquad
r = \frac{\mathrm{Cov}(x,y)}{\sigma_x \sigma_y}.
\]

\begin{figure}[H]
\centering
\begin{tikzpicture}
\begin{axis}[
  width=11cm, height=7cm,
  xlabel={\( x \)},
  ylabel={\( y \)},
  grid=major,
  grid style={dashed,gray!40},
  xmin=0, xmax=10,
  ymin=0, ymax=20,
  legend pos=north west,
  legend style={font=\footnotesize},
]
\addplot[only marks, mark=*, mark size=2pt, blue]
  coordinates {
   (1,3.2) (2,5.1) (3,6.8) (4,9.4) (5,10.1)
   (6,12.9) (7,13.8) (8,16.2) (9,17.1) (10,18.9)
  };
\addplot[red, thick, domain=0:10, samples=2] {1.78*x + 1.28};
\legend{Points, Droite \( y = 1{,}78x + 1{,}28 \)}
\end{axis}
\end{tikzpicture}
\caption{Exemple de nuage de points et droite de régression linéaire.}
\end{figure}

\begin{calcsheet}
73 & Pente \( a \) & \fc{=PENTE(B2:B11;A2:A11)} & 1,78 & \\ \hline
74 & Ordonnée \( b \) & \fc{=ORDONNEE.ORIGINE(B2:B11;A2:A11)} & 1,28 & \\ \hline
75 & \( r \) & \fc{=COEFFICIENT.CORRELATION(A2:A11;B2:B11)} & 0,998 & \\ \hline
76 & \( r^2 \) & \fc{=COEFFICIENT.DETERMINATION(B2:B11;A2:A11)} & 0,996 & \\ \hline
77 & \( a' = 1/a \) & \fc{=1/D73} & 0,562 & \\ \hline
78 & \( b' = -b/a \) & \fc{=-D74/D73} & -0,719 & \\ \hline
\end{calcsheet}

% ============================================================
\section{Comparaison visuelle : exact vs.\ classes}

\begin{figure}[H]
\centering
\begin{tikzpicture}
\begin{axis}[
  width=12cm, height=6.5cm,
  ybar,
  bar width=1.2cm,
  xlabel={Indicateur},
  ylabel={Valeur},
  symbolic x coords={Moyenne, Variance, Ecart-type},
  xtick=data,
  ymin=0, ymax=17,
  legend pos=north east,
  legend style={font=\footnotesize},
  nodes near coords,
  nodes near coords style={font=\scriptsize},
]
\addplot coordinates {(Moyenne,15.835) (Variance,3.032) (Ecart-type,1.741)};
\addplot coordinates {(Moyenne,16.0) (Variance,2.8) (Ecart-type,1.673)};
\legend{Exact, Classes}
\end{axis}
\end{tikzpicture}
\caption{Comparaison des indicateurs par calcul exact et par classes.}
\end{figure}

% ============================================================
\section{Modus operandi — Générer les graphiques avec LibreOffice Calc}

\subsection{Préparation de la feuille}

\begin{enumerate}[label=\textbf{\arabic*.}]
  \item Ouvrir LibreOffice Calc.
  \item Saisir les 40 valeurs dans \( A2:A41 \).
  \item Ajouter les formules auxiliaires (moyenne, médiane, quartiles,
        écarts-types) comme indiqué dans les simulations précédentes.
\end{enumerate}

\subsection{Histogramme des effectifs}

\begin{enumerate}[label=\textbf{\arabic*.}]
  \item Construire le tableau des classes dans deux colonnes :
        \( B2:B5 \) = centres (ou bornes), \( C2:C5 \) = effectifs.
  \item Sélectionner \( B2:C5 \).
  \item Menu \textsf{Insertion \(\to\) Graphique}.
  \item Type : \textsf{Colonnes} (ou \textsf{Barres}).
  \item Cliquer \textsf{Suivant}, puis \textsf{Terminer}.
  \item Pour transformer en histogramme : clic droit sur une barre
        \(\to\) \textsf{Écart} \(\to\) mettre \SI{0}{\percent}
        (barres jointives).
  \item Double-clic sur l'axe X \(\to\) \textsf{Étiquettes} \(\to\)
        choisir les bornes de classes.
\end{enumerate}

\subsection{Boîte à moustaches (à partir de LibreOffice 7.1)}

\begin{enumerate}[label=\textbf{\arabic*.}]
  \item Sélectionner la plage \( A2:A41 \) (données brutes).
  \item \textsf{Insertion \(\to\) Graphique}.
  \item Type : \textsf{Boîte à moustaches}.
  \item Suivant / Terminer.
  \item Pour afficher les valeurs : double-clic sur la boîte
        \(\to\) onglet \textsf{Étiquettes de données} \(\to\) cocher
        \textsf{Afficher les valeurs}.
  \item Le seuil de Tukey est calculé automatiquement ;
        pour le modifier : clic droit \(\to\) \textsf{Format de la série}
        \(\to\) \textsf{Options de la boîte} \(\to\) coefficient
        d'écart interquartile (défaut : 1,5).
\end{enumerate}

\subsection{Nuage de points et droite de régression}

\begin{enumerate}[label=\textbf{\arabic*.}]
  \item Saisir \( x \) dans \( A2:A11 \) et \( y \) dans \( B2:B11 \).
  \item Sélectionner \( A2:B11 \).
  \item \textsf{Insertion \(\to\) Graphique} \(\to\) type
        \textsf{XY (dispersion)}.
  \item Choisir \textsf{Points seuls} puis \textsf{Terminer}.
  \item Clic droit sur un point \(\to\) \textsf{Insérer une courbe de tendance}.
  \item Choisir \textsf{Linéaire}, cocher
        \textsf{Afficher l'équation} et \textsf{Afficher \( R^2 \)}.
  \item Pour la droite inverse : échanger les colonnes en abscisse
        et en ordonnée, puis répéter.
\end{enumerate}

\subsection{Coefficient de variation et tableau récapitulatif}

\begin{enumerate}[label=\textbf{\arabic*.}]
  \item Ajouter une cellule \( D60 \) : \fc{=ECARTYPE.P(A2:A41)/MOYENNE(A2:A41)}.
  \item Format \(\to\) \textsf{Nombre} \(\to\) \textsf{Pourcentage}.
  \item Pour un récapitulatif graphique : sélectionner les
        cellules d'indicateurs et \textsf{Insertion \(\to\) Graphique
        \(\to\) Colonnes}.
\end{enumerate}

% ============================================================
\section{Récapitulatif général}

\begin{center}
\small
\begin{tabular}{|l|c|c|}
\hline
\rowcolor{calchead}
\textbf{Indicateur} & \textbf{Valeur exacte} & \textbf{Valeur par classes} \\ \hline
Effectif \( n \)            & 40      & 40 \\ \hline
\( \sum x_i \)              & 633,4   & --- \\ \hline
\( \sum x_i^2 \)            & 10\,151,16 & --- \\ \hline
Moyenne \( \bar{x} \)       & 15,835  & 16,0 \\ \hline
Médiane \( Me \)            & 15,85   & --- \\ \hline
\( Q_1 \)                   & 14,55   & --- \\ \hline
\( Q_3 \)                   & 17,15   & --- \\ \hline
\( IQR \)                   & 2,60    & --- \\ \hline
Étendue                     & 7,0     & --- \\ \hline
Variance \( V \) (pop.)     & 3,032   & 2,8 \\ \hline
Écart-type \( \sigma \)     & 1,741   & 1,673 \\ \hline
Écart-type \( s \) (éch.)   & 1,763   & --- \\ \hline
\( CV \)                    & 11,0\,\% & 10,5\,\% \\ \hline
\end{tabular}
\end{center}

\section*{Conclusion biologique}

Le lot présente une concentration moyenne de
\( 15{,}835\ \si{g/L} \), très proche de la cible \( 15{,}5\ \si{g/L} \),
avec un \( CV \) d'environ \( 11\ \% \) indiquant une bonne homogénéité.
Aucune valeur aberrante n'est détectée selon la règle de Tukey
(borne supérieure : \( 21{,}05\ \si{g/L} \) ; max observé :
\( 19{,}1\ \si{g/L} \)). Le lot peut être déclaré
\textbf{conforme}.

\end{document}